% HD-G30. GPT-6 Astra Pro /root/g34; inherited effort; 21 November 2026.
\section{HD-G30: the actual matched-pair reflection factors}
\label{sec:matched-pair-reflection-factors}

This section resolves the matched-pair case in the actual
normal-subgroup decomposition used in HD-S42. Let \(G\) be
the finite absolutely irreducible coefficient group on \(V_K=K^m\) over the
imaginary quadratic exact character field \(K\). Let
\(N\triangleleft G\) be generated by a \(G\)-stable set
\(\mathcal X\) of actual rank-two elements. In the case at
issue the restriction to \(N\) is multiplicity-free, its
complex irreducible constituents \(W_i\) all have dimension
\(r\), and the supports
\[
 S_x=\{i:(\rho(x)-I)|_{W_i}\ne0\},\qquad x\in\mathcal X,
                                                               \tag{G30.1}
\]
are the edges of a perfect matching, with possible repetitions.
The constituent projectors are not individually defined over
\(K\). These are the precise data of this case; no stability
of an arbitrary complex system of imprimitivity is assumed.

For every \(r\ge8\), the resulting natural primitive
sixfold period map has rank at most seven. For \(r\le7\),
it is the period map of an actual coefficient group of rank
\(2r\le14\), with the polarized quotient maps below.
Consequently, over \(K=\mathbb Q(i)\) this case gives only
the already established ordinary determinants \([-1],[-3]\)
at period rank nine, and over
\(K=\mathbb Q(\sqrt{-d})\),
\(d\in\{2,7,11,19,43,67,163\}\), it gives only
\([-1]\), by HD-S40. The high-rank bound itself holds over
every imaginary quadratic field.

\paragraph{1. The quadratic projector field and its actual action.}
Put \(M=K(\zeta_h)\), where \(h\) is the exponent of
\(N\). For each constituent, let \(\chi_i\) be its
character and form the actual matrix
\[
 p_i=\frac r{|N|}\sum_{n\in N}\chi_i(n^{-1})\rho(n).
                                                               \tag{G30.2}
\]
Character averaging gives \(p_ip_j=\delta_{ij}p_i\),
\(\sum_i p_i=I\), and image \(W_i\). More explicitly,
averaging \(A\mapsto\rho_j(n)A\rho_i(n^{-1})\) is the
projection onto \(\operatorname{Hom}_N(W_i,W_j)\), whose
dimension is \(\delta_{ij}\). Its trace proves the
displayed scalar action of (G30.2) on every constituent.
The same averaging with individual matrix coefficients gives
all matrix units on each constituent, exactly as in
(G29.7), with \(2\) replaced by \(r\). Thus these are
precisely the primitive central idempotents of the actual
complexified image algebra, each occurring once.

The Galois action on these projectors commutes with conjugation
by every original \(K\)-matrix \(\rho(g)\). The action
of \(G\) on the labels is transitive: the sum of the
constituents in an orbit would otherwise be a proper
\(G\)-invariant complex subspace. A permutation commuting
with a transitive group and fixing one label fixes all labels,
by carrying that label through its whole orbit. Also Galois
fixes each individual set \(S_x\), since
\[
 \gamma((\rho(x)-I)p_i)=(\rho(x)-I)p_{\gamma i}.
                                                               \tag{G30.3}
\]
Hence every nonidentity Galois permutation swaps the two
endpoints of every matching edge. There is exactly one such
permutation. Its kernel has index two, and its fixed field
is an actual quadratic extension \(E/K\) inside \(M\).
Write \(\sigma\) for its nontrivial automorphism.
All \(p_i\) lie in \(M_m(E)\), and on each pair
\(\{i,i'\}\) one has \(p_{i'}=\sigma(p_i)\).
Thus
\[
 e=p_i+p_{i'}\in M_m(K),\qquad U=eV_K,\qquad\dim_KU=2r.
                                                               \tag{G30.4}
\]
The original algebra base change is literal: a basis of the
\(K\)-matrix span of \(\rho(N)\) remains independent
over \(M\), by a nonzero maximal minor. Its Galois fixed
part is therefore that original matrix span.
The pair subspaces are orthogonal for the prescribed positive
invariant coefficient form. Indeed character unitarity gives
\(\chi_i(n^{-1})=\overline{\chi_i(n)}\), and taking the
Hermitian adjoint of (G30.2), followed by \(n\mapsto n^{-1}\),
gives \(p_i^\dagger=p_i\). The mutually orthogonal
idempotents and their pair sums are therefore the actual
orthogonal projectors for this form.

There is an explicit field embedding on this actual subspace:
\[
 \iota:E\longrightarrow\operatorname{End}_K(U),\qquad
             a\longmapsto a p_i+\sigma(a)p_{i'}.
                                                               \tag{G30.5}
\]
The matrix is fixed by \(\operatorname{Gal}(M/K)\),
is in the image algebra just described, and its products
are exactly \(\iota(a)\iota(b)=\iota(ab)\).
It turns \(U\) into an \(r\)-dimensional \(E\)-space;
write \(U_E\) for \(U\) equipped with this action.
The actual comparison with the first constituent is
\[
 U_E\longrightarrow p_iE^m,\quad u\longmapsto p_i u,
 \qquad w\longmapsto w+\sigma(w)\quad\hbox{as inverse}.
                                                               \tag{G30.6}
\]
Both identities follow by applying the two orthogonal
idempotents. The inverse takes values in \(K^m\), and
(G30.5) makes these maps \(E\)-linear.

Each generator supported on this pair is the identity on
all other pairs. The subgroups generated by the generators
on distinct pairs commute and have trivial intersections,
because their matrices act on disjoint direct summands.
Consequently \(N\) is their actual direct product. Denote
its factor on \(U\) by \(L\). It acts faithfully and
absolutely irreducibly on the \(E\)-space in (G30.6):
faithfulness follows by applying \(\sigma\) to a matrix
acting trivially on its first constituent, and irreducibility
is the defining irreducibility of \(W_i\). Each of its
specified generators has complex rank one on each of the two
constituents, since its total rank is two. Hence \(L\)
is an actual finite \(E\)-linear reflection group of rank
\(r\).

Let \(F=\mathbb Q(\operatorname{tr}_E L)\subset E\).
Then
\[
                              KF=E.                         \tag{G30.7}
\]
Indeed \(KF=K\) would make the two constituent characters
equal under \(\sigma\), hence equal as characters of
\(N\), contrary to their inequivalence. The only other
intermediate field of \(E/K\) is \(E\) itself.
The stabilizer \(H\) of the pair in \(G\) normalizes
\(L\), and conjugates \(\iota(a)\) to \(\iota(a)\)
or \(\iota(\sigma(a))\) according as it preserves or
interchanges the two constituents. Thus its action on \(U\)
is literally \(E/K\)-semilinear. Moreover some member
interchanges them: transitivity of \(G\) on all constituent
labels supplies a member carrying one endpoint to its partner,
and it then stabilizes this edge.

\paragraph{2. The high-rank reflection list and an exact trace recovery.}
For \(r\ge8\), Cohen's classification gives either a
real rational reflection type or an imprimitive group
\(G(\ell,e,r)\), \(e\mid\ell\). The precise inputs
are Cohen, \emph{Finite complex reflection groups},
Ann. Sci. ENS (4) \textbf{9} (1976), 379--436:
Theorem (2.4), pp.386--387; the real list (4.4)(i), p.399;
Theorem (5.12), p.435 and table (4.16), p.412.
The nonreal primitive list has maximum rank six. The real
types of rank at least eight have rational matrices: type
\(A_r\) acts on the sum-zero subspace of \(\mathbb Q^{r+1}\),
types \(B_r,D_r\) by signed permutations, and the additional
rank-eight type \(E_8\) by its integral Cartan reflections.
Their characters are rational. This contradicts (G30.7),
so the actual \(L\) is of the imprimitive type.
The 1978 erratum, printed p.613, changes (5.14) and deletes
(5.15); neither is used here.

We now recover the field required to define its individual
axes, without assuming that a complex conjugating matrix is
already defined over \(E\). Fix the matrix \(T\) supplied
by this classification, so that
\(T\rho_E(L)T^{-1}=G(\ell,e,r)\). Every standard
matrix \(d\) used below therefore corresponds to the
literal actual matrix \(T^{-1}dT\in\rho_E(L)\).
In that model choose a primitive \(\ell\)-th root \(\zeta\).
Irreducibility gives \(\ell>1\). Both actual group
elements corresponding to
\[
 \operatorname{diag}(\zeta,\zeta^{-1},1,\ldots,1),\qquad
 \operatorname{diag}(\zeta,\zeta,\zeta^{-2},1,\ldots,1)
                                                               \tag{G30.8}
\]
belong to every \(G(\ell,e,r)\). Subtracting the indicated
identity entries from their traces gives
\[
 c=\zeta+\zeta^{-1},\qquad v=2\zeta+\zeta^{-2},\qquad
           \zeta=\frac{v-c^2+1}{2-c}.                     \tag{G30.9}
\]
The denominator is nonzero, since a root of unity with
\(\zeta+\zeta^{-1}=2\) equals one. Conversely all entries
in the standard monomial model are in \(\mathbb Q(\zeta)\).
Trace is unchanged under its complex conjugacy with \(L\),
so these two inclusions prove
\[
             F=\mathbb Q(\zeta),\qquad E=K(\zeta).
                                                               \tag{G30.10}
\]
In particular \(\sigma(\zeta)\ne\zeta\).

\paragraph{3. Characteristic diagonal kernel and actual rank-two projectors.}
Let \(D\triangleleft L\) be the diagonal subgroup in
the standard model. It is intrinsic to the abstract group:
it is the unique largest abelian normal subgroup. Indeed
\(L/D\simeq S_r\), and an abelian normal subgroup has
abelian normal image in \(S_r\). For \(r\ge5\) such
an image is trivial: its intersection with the simple
\(A_r\) is trivial, so its order is at most two; a normal
subgroup of order two is central, while \(S_r\) has trivial
centre. Thus every abelian normal subgroup lies in \(D\).
In particular every semilinear normalizer of \(L\)
normalizes this actual \(D\).

The coordinate characters \(\lambda_j:D\to\mu_\ell\)
are pairwise distinct. For \(i\ne j\), choose
\(k\ne i,j\); the element with entries \(\zeta\)
at \(i\), \(\zeta^{-1}\) at \(k\), and one elsewhere
distinguishes \(\lambda_i\) from \(\lambda_j\).
By (G30.10) every character value lies in \(E\). Hence
the matrices
\[
 f_j=\frac1{|D|}\sum_{t\in D}\lambda_j(t^{-1})\rho_E(t)
                         \in\operatorname{End}_E(U_E)
                                                               \tag{G30.11}
\]
are actual rank-one \(E\)-projectors. Character averaging
on the abelian group proves that they are mutually orthogonal,
sum to one, and have exactly the corresponding complex axes
as their scalar extensions. Thus their images themselves
supply an \(E\)-basis realizing the standard axes.

On the original \(K\)-space their literal matrices are
\[
 q_j=\frac1{|D|}\sum_{t\in D}
           \iota(\lambda_j(t^{-1}))\rho(t)e,
 \qquad U_j=q_jV_K,\qquad\dim_KU_j=2.                    \tag{G30.12}
\]
These are the underlying \(K\)-linear maps of (G30.11),
not projectors on a separately chosen complex representation.
Every semilinear normalizer permutes their images, since it
normalizes \(D\) and therefore permutes its simultaneous
one-dimensional \(E\)-eigenspaces. Internal permutations
in \(L\) act transitively on these axes. Carrying the
construction to all pairs gives a transitive actual
\(K\)-rank-two block system for the full original \(G\).

The prescribed positive invariant Hermitian coefficient form
makes these \(K\)-blocks orthogonal. For different axes
in the same pair use the element just used to distinguish
their characters. It acts identically on the second axis,
while \(\rho(t)-I\) is invertible on the first: its two
complex eigenvalues there are \(\zeta-1\) and
\(\sigma(\zeta)-1\), both nonzero. Invariance of the
form therefore gives zero cross pairing. For axes in distinct
pairs the same element is identity on the entire other pair.
This proves orthogonality for the original coefficient form.

The stabilizer \(J\) of one \(K\)-axis has an absolutely
irreducible two-dimensional representation \(\tau\).
Its two complex \(D\)-characters are \(\lambda_j\)
and \(\sigma\lambda_j\), distinct by (G30.10).
An element of \(H\) swapping the two constituent types
may send this axis to another. Compose it with an internal
permutation in \(L\) carrying that axis back. The resulting
element lies in \(J\) and swaps these two eigenspaces.
Every \(D\)-invariant complex subspace is a sum of the
two distinct character lines, and this interchange excludes
either line as a \(J\)-invariant subspace.

For actual left coset representatives \(t_a\), the map
\[
       \bigoplus_a U_j\longrightarrow V_K,
                     (u_a)\longmapsto\sum_a\rho(t_a)u_a
                                                               \tag{G30.13}
\]
is the explicit induction isomorphism, with the actual
\(\tau\)-transition when \(gt_a=t_bh\).
All traces of \(\tau\) lie in \(K\). If they were
rational, the induced trace formula from (G30.13) would
make the original character rational. Thus its exact
character field is the original \(K\).

\paragraph{4. Actual covers, primitive corners, and the rank bound.}
Write \(J_0=\ker\tau\). From the original full cover
\(\widetilde C\to\mathbb P^1\), take the actual quotients
\[
       Y=\widetilde C/J,\qquad C'=\widetilde C/J_0,
                    \alpha:\widetilde C\to C'.             \tag{G30.14}
\]
For the block projection \(\pi=q_j\), the coefficient
sheaf maps and inverse are
\[
 \begin{aligned}
 F_v(x)&=\pi\rho(x)v,& F_v(hx)&=\tau(h)F_v(x),\\
 F_{\rho(g)v}(x)&=F_v(xg),&
 F&\longmapsto\sum_{x\in J\backslash G}\rho(x)^{-1}F(x).
 \end{aligned}                                              \tag{G30.15}
\]
The left-sheet label \(Jt_a^{-1}\) is the block \(t_aJ\).
On each local branch orbit both sides have the fixed space
of the same product monodromy. Thus these are actual
middle-extension maps, including the branch stalks.

Let \(m=\dim_K V\), choose \(0\ne v\in U_j\), and
put \(q=h(v,v)\). Orthogonality makes every coefficient
outside \(J\) vanish, while \(m/|G|=2/|J|\). Therefore
the actual rational matrix-coefficient corner idempotents obey
\[
 e_G=\frac{2m}{|G|q}\sum_{g\in G}
       \operatorname{Re}h(v,\rho(g^{-1})v)g
     =\frac4{|J|q}\sum_{h\in J}
       \operatorname{Re}h(v,\tau(h^{-1})v)h=e_J.
                                                               \tag{G30.16}
\]
The same integral multiple has the same abelian image inside
\(J\widetilde C\). The quotient pullback carries the
corresponding primitive factor of \(JC'\) isogenously onto
it and satisfies
\[
 \alpha_*\alpha^*=[|J_0|],\qquad
   (\alpha^*)^*\theta_{J\widetilde C}=|J_0|\theta_{JC'}.
                                                               \tag{G30.17}
\]
It commutes with the actual quotient deck operators and hence
with the same primitive-corner field action. The two Schur
constants agree: \(2|G|/(mq)=2|J|/(2q)\).

An originally specified primitive line \(Kw\) is retained
by the actual matrix coefficient of
\(T_{v,w}(x)=v h(w,x)/h(w,w)\):
\[
 b=\frac{2m}{|G|}\sum_{g\in G}
        \operatorname{Re}\operatorname{tr}
                (T_{v,w}\rho(g^{-1}))g,\qquad
 b^\dagger b=\frac q{h(w,w)}e_w.                            \tag{G30.18}
\]
A positive integer \(M_0\) clearing its coefficients gives
the actual isogeny between the two primitive abelian factors,
with exact pullback polarization multiplier
\(M_0^2q/h(w,w)\). These are the literal averaging and
Rosati maps of HD-G29, now on the descended axis. Thus the
original period map factors through the rank-two quotient
family after fixing the finite monodromy labels; it is not
replaced merely by a dimensionally similar Hodge structure.

The faithful rank-two coefficient sheaf on \(Y\) has
parabolic cohomology of \(K\)-dimension six. With
\(g=g(Y)\) and \(b\) nontrivial local coefficient
branches, irreducibility and dual irreducibility give
\[
                 6=2(2g-2)+\sum_{a=1}^b c_a,
                                      \qquad c_a\ge1.       \tag{G30.19}
\]
For \(g=0\) there are at most ten branches, so the
pointed-base dimension is at most seven. For \(g=1\)
there are at most six branches and dimension at most six.
For \(g=2\) there are at most two branches and dimension
at most five. Larger genus is impossible. The finite cover
choices over a fixed pointed base are discrete, so these
bound the actual primitive period map as well. This proves
the high-rank conclusion.

\paragraph{5. The lower-rank pairs and the retained determinant.}
For \(r\le7\), use the entire actual pair \(U\),
not an unproved axis system. Its stabilizer \(H\) acts
absolutely irreducibly on its \(K\)-space of dimension
\(2r\): the two inequivalent irreducible \(N\)-types
are interchanged, as proved in paragraph1. The induced
character again gives exact character field \(K\).
The same quotient, coefficient, corner and polarization maps
(G30.14)--(G30.18), with \(J\) replaced by \(H\)
and coefficient rank two replaced by \(2r\), retain the
original primitive sixfold family.

Rank \(2r=2\) has the preceding rank-seven bound. For
\(2r\ge4\), the formula
\[
                   6=2r(2g(Y)-2)+\sum c_a                 \tag{G30.20}
\]
excludes \(g(Y)\ge2\), while genus one gives at most
six branches and period rank at most six. Therefore a
rank-nine period family has \(Y=\mathbb P^1\). Its
faithful coefficient group has rank \(2r\le14\), so
HD-S40 applies to this actual source, including its stated
seven-field extension. Every polarization multiplier above
is positive rational; in Hermitian dimension six its
determinant multiplier is \(c^6=N_{K/\mathbb Q}(c^3)\).
Hence the resulting determinant conclusions are exactly those
of the originally specified primitive factor.

The proof closes this actual matched-pair case. It does not
assert that a general normal subgroup generated by
bireflections has multiplicity one, or that it is nontrivial
in every larger coefficient rank; those are separate cases
in the normal-subgroup analysis.
